% GENERATED FILE. Edit canonical lessons, then run npm run generate:books.
% canonical-source-hash: fnv1a64:de950180e1d53ddc
% canonical-lessons: UR-C39-pet, UR-C39-kandha, UR-C39-ungli, UR-C39-hot, UR-C39-gardan

\chapter{Stomach, Shoulder, Finger, Lip, Neck}
\label{ch:ur-stomach-shoulder-finger-lip-neck}
\hlchaptermodality{39}

\begin{chapteropening}
\textbf{By the end of this chapter:} \emph{I can say the stomach, a shoulder, a finger, a lip and the neck.}

Complete the last lesson of chapter 39: I can say the stomach, a shoulder, a finger, a lip and the neck.
\end{chapteropening}

\section[\texorpdfstring{\ur{پیٹ} (peṭ)}{پیٹ (peṭ)}]{\ur{پیٹ} (peṭ) — the stomach}

\label{lesson:UR-C39-pet}



\begin{quote}
Before the new one: say the Urdu for a tooth, then the Urdu for the throat.
\end{quote}

\subsection*{You'll want to know: \ur{پیٹ}}
\textbf{\ur{پیٹ}} — \emph{peṭ} — \textquotedblleft{}the stomach\textquotedblright{}.

\begin{cousinweb}
\textbf{\ur{پیٹ}} (\emph{peṭ}) is inherited. CDIAL lists it under a reconstructed \emph{*peṭṭa}, \textquotedblleft{}belly\textquotedblright{}, and notes that it may be the same word as similar forms meaning \textquotedblleft{}basket\textquotedblright{} (CDIAL 8376). DEDR compares its variant \emph{poṭṭa} with Dravidian words for \textquotedblleft{}belly\textquotedblright{}, such as Telugu \emph{poṭṭa} (DEDR 4494).
\end{cousinweb}

\begin{grammarlens}[title={What you've built}]
One more word for this chapter.
\end{grammarlens}

\subsection*{Your turn}
\begin{itemize}
\raggedright
  \item \emph{Say it:} \emph{peṭ}
  \item \emph{Say it:} \emph{peṭ}, once more
  \item \emph{Say it:} say \emph{peṭ}
  \item \emph{Recall:} say the Urdu for a tooth, then the Urdu for the throat, then say \emph{peṭ} again
\end{itemize}

\begin{tcolorbox}[breakable,colback=teal!4,colframe=teal!35!black,title={Before you move on}]
What does \ur{پیٹ} mean? (\textquotedblleft{}The stomach\textquotedblright{}.) Say it once more.
\end{tcolorbox}
\section[\texorpdfstring{\ur{کندھا} (kandhā)}{کندھا (kandhā)}]{\ur{کندھا} (kandhā) — a shoulder}

\label{lesson:UR-C39-kandha}



\begin{quote}
Before the new one: say the Urdu for the throat, then the Urdu for the stomach.
\end{quote}

\subsection*{You'll want to know: \ur{کندھا}}
\textbf{\ur{کندھا}} — \emph{kandhā} — \textquotedblleft{}a shoulder\textquotedblright{}.

\begin{cousinweb}
\textbf{\ur{کندھا}} (\emph{kandhā}) is a tadbhava, a word inherited from Sanskrit and reshaped by sound change. It comes from \emph{skandha}, \textquotedblleft{}shoulder\textquotedblright{} (CDIAL 13627). CDIAL allows that a form without the \emph{s-}, \emph{*kandha}, existed beside it.
\end{cousinweb}

\begin{grammarlens}[title={What you've built}]
One more word for this chapter.
\end{grammarlens}

\subsection*{Your turn}
\begin{itemize}
\raggedright
  \item \emph{Say it:} \emph{kandhā}
  \item \emph{Say it:} \emph{kandhā}, once more
  \item \emph{Say it:} say \emph{kandhā}
  \item \emph{Recall:} say the Urdu for the throat, then the Urdu for the stomach, then say \emph{kandhā} again
\end{itemize}

\begin{tcolorbox}[breakable,colback=teal!4,colframe=teal!35!black,title={Before you move on}]
What does \ur{کندھا} mean? (\textquotedblleft{}A shoulder\textquotedblright{}.) Say it once more.
\end{tcolorbox}
\section[\texorpdfstring{\ur{انگلی} (unglī)}{انگلی (unglī)}]{\ur{انگلی} (unglī) — a finger}

\label{lesson:UR-C39-ungli}



\begin{quote}
Before the new one: say the Urdu for the stomach, then the Urdu for a shoulder.
\end{quote}

\subsection*{You'll want to know: \ur{انگلی}}
\textbf{\ur{انگلی}} — \emph{unglī} — \textquotedblleft{}a finger\textquotedblright{}.

\begin{cousinweb}
\textbf{\ur{انگلی}} (\emph{unglī}) is a tadbhava, a word inherited from Sanskrit and reshaped by sound change. It comes from \emph{aṅguli}, \textquotedblleft{}finger, toe\textquotedblright{} (CDIAL 135).
\end{cousinweb}

\begin{grammarlens}[title={What you've built}]
One more word for this chapter.
\end{grammarlens}

\subsection*{Your turn}
\begin{itemize}
\raggedright
  \item \emph{Say it:} \emph{unglī}
  \item \emph{Say it:} \emph{unglī}, once more
  \item \emph{Say it:} say \emph{unglī}
  \item \emph{Recall:} say the Urdu for the stomach, then the Urdu for a shoulder, then say \emph{unglī} again
\end{itemize}

\begin{tcolorbox}[breakable,colback=teal!4,colframe=teal!35!black,title={Before you move on}]
What does \ur{انگلی} mean? (\textquotedblleft{}A finger\textquotedblright{}.) Say it once more.
\end{tcolorbox}
\section[\texorpdfstring{\ur{ہونٹ} (hõṭ)}{ہونٹ (hõṭ)}]{\ur{ہونٹ} (hõṭ) — a lip}

\label{lesson:UR-C39-hot}



\begin{quote}
Before the new one: say the Urdu for a shoulder, then the Urdu for a finger.
\end{quote}

\subsection*{You'll want to know: \ur{ہونٹ}}
\textbf{\ur{ہونٹ}} — \emph{hõṭ} — \textquotedblleft{}a lip\textquotedblright{}.

\begin{cousinweb}
\textbf{\ur{ہونٹ}} (\emph{hõṭ}) is inherited. CDIAL lists Hindi \emph{hõṭh}, \textquotedblleft{}lip\textquotedblright{}, under Sanskrit \emph{oṣṭha}, \textquotedblleft{}lip\textquotedblright{} (CDIAL 2563); its addenda also list \emph{hõṭh} under a reconstructed compound \emph{*adhamauṣṭha}, \textquotedblleft{}lower lip\textquotedblright{}, which would account for the \emph{h-} (CDIAL 247a).
\end{cousinweb}

\begin{grammarlens}[title={What you've built}]
One more word for this chapter.
\end{grammarlens}

\subsection*{Your turn}
\begin{itemize}
\raggedright
  \item \emph{Say it:} \emph{hõṭ}
  \item \emph{Say it:} \emph{hõṭ}, once more
  \item \emph{Say it:} say \emph{hõṭ}
  \item \emph{Recall:} say the Urdu for a shoulder, then the Urdu for a finger, then say \emph{hõṭ} again
\end{itemize}

\begin{tcolorbox}[breakable,colback=teal!4,colframe=teal!35!black,title={Before you move on}]
What does \ur{ہونٹ} mean? (\textquotedblleft{}A lip\textquotedblright{}.) Say it once more.
\end{tcolorbox}
\section[\texorpdfstring{\ur{گردن} (gardan)}{گردن (gardan)}]{\ur{گردن} (gardan) — the neck}

\label{lesson:UR-C39-gardan}



\begin{quote}
Before the new one: say the Urdu for a finger, then the Urdu for a lip.
\end{quote}

\subsection*{You'll want to know: \ur{گردن}}
\textbf{\ur{گردن}} — \emph{gardan} — \textquotedblleft{}the neck\textquotedblright{}.

\begin{cousinweb}
\textbf{\ur{گردن}} (\emph{gardan}) is a loan, usually traced to Persian \emph{gardan}, \textquotedblleft{}neck\textquotedblright{}.
\end{cousinweb}

\begin{grammarlens}[title={What you've built}]
One more word for this chapter.
\end{grammarlens}

\subsection*{Your turn}
\begin{itemize}
\raggedright
  \item \emph{Say it:} \emph{gardan}
  \item \emph{Say it:} \emph{gardan}, once more
  \item \emph{Say it:} say \emph{gardan}
  \item \emph{Recall:} say the Urdu for a finger, then the Urdu for a lip, then say \emph{gardan} again
\end{itemize}

\begin{tcolorbox}[breakable,colback=teal!4,colframe=teal!35!black,title={Before you move on}]
What does \ur{گردن} mean? (\textquotedblleft{}The neck\textquotedblright{}.) Say it once more.
\end{tcolorbox}
